% V10: connected journal-first research argument and selective mathematical notation.
\newcommand{\HalfMoonIntroduction}{%
\section{Introduction}\label{sec:introduction}
Normalizing observations to unit length suppresses differences in scale by retaining direction and discarding distance from a chosen origin. If that distance helps distinguish classes, observations requiring different labels can become identical before a classifier sees them. We investigate whether choosing a different origin can prevent this loss and how little distance information must be restored when it cannot. Geometric proofs and controlled neural experiments connect the retained information to the smallest achievable classification error.

The half-moon problem makes this difficulty visible. Each observation is a point labeled by which of two crescent-shaped regions generated it (Fig.~\ref{fig:geometry}). Two points on the same ray share a direction even when their labels differ. Moving the origin changes which points merge, making centering part of the information-preservation problem. We measure the effect through Bayes risk, the smallest population classification error attainable from a given representation. This separates information discarded by preprocessing from additional error introduced by learning.

Recent \emph{Pattern Recognition Letters} research approaches this distinction from complementary directions. Poblete et al.\ \citep{Poblete2026} hold learned representations fixed and show that classifier rankings vary with the dataset. Ahmadian and Lindsten\ \citep{Ahmadian2025} observe poorer novelty-detection performance after unit normalization and suggest that lost radial information contributes. These findings motivate asking what the representation itself makes possible. Decision-sufficiency and lossless-transformation theory already characterize preservation of optimal decisions \citep{Silva2012,Silva2024,Gyorfi2023}; the geometric question is how that preservation behaves as the normalization origin changes over the whole plane.

A positive error at each tested origin does not establish a positive lower bound across all origins: increasingly distant centers might still make the loss arbitrarily small. We resolve this issue by proving a uniform limit to orthogonal-projection risk as origins recede. Connected planar regions favoring opposite decisions then yield a global loss criterion and computable lower certificates. For a specified half-moon family, these results establish unavoidable loss in both disjoint and overlapping configurations. In one disjoint midpoint configuration, a single near/far radius measurement selects between two angular rules and restores every class decision. Its exact operating interval determines the permissible radius-measurement error. Paired neural studies compare learned errors with these representation-specific limits.
}

\newcommand{\HalfMoonBackground}{%
\section{Background and related work}\label{sec:background}
\paragraph{What direction retains.}
The spatial-sign transformation retains the direction of a centered observation \citep{Serneels2006}. For an observation $x$ and shared origin $c$, it is $T_c(x)=(x-c)/\|x-c\|$ when $x\ne c$. For illustration, $(1,1)$ and $(2,2)$ become identical about $(0,0)$; about $(0,1)$ their directions differ. Keeping the radius $\rho=\|x-c\|$ alongside direction $u$ permits reconstruction through $x=c+\rho u$. This gives a full-information control against which to assess discarding radius.

\paragraph{Geometric evidence in recent classification research.}
Poblete et al.\ \citep{Poblete2026} compare four classifiers on the same fixed feature vectors from trained generative models across eight datasets. Their class-centered radial-compatibility descriptor measures direction-independent spread within classes and radial separation between classes. Classifier rankings vary, and the predicted directional association between this descriptor and classifier advantage is unsupported. This separates a useful description of geometry from a guaranteed decision rule. Our question instead concerns how moving a shared normalization origin changes the observations that become indistinguishable.

Ahmadian and Lindsten\ \citep{Ahmadian2025} combine radial scoring with context quantization for human-activity novelty detection. Their separate normalization and quantization ablations show that the normalized variant performs worse than the unnormalized radial-kernel variant without quantization on both principal datasets. They suggest information loss under normalization as an explanation. Yang et al.\ \citep{Yang2026Compression} examine another consequence of compression: image-classification predictions can remain correct while visual explanations deteriorate, motivating an explanation-aware training objective. Together, these studies show why retained task information must be specified explicitly. Here the object is the optimal class decision after a fixed direction map, before classifier fitting.

\paragraph{Preserved decisions and restricted radial recovery.}
The Bayes decision selects the more probable class at an input. Silva and Narayanan\ \citep{Silva2012} analyze representation-induced error; Silva et al.\ \citep{Silva2024} distinguish retaining information about the label from retaining the optimal decision. Gy\"orfi et al.\ \citep{Gyorfi2023} characterize losslessness across loss functions through conditional independence and derive excess-risk upper bounds. Encoder--decoder analysis also relates information loss to prediction losses \citep{Silva2025}. These results establish the statistical target; a direction map requires a geometric account of which opposing decisions it merges.

Dubois et al.\ \citep{Dubois2021} characterize compression for invariant predictive tasks and identify unit direction as a maximal invariant for scaling: it gives the same encoding precisely to points on the same ray. Their task condition requires the conditional label law to be constant within each equivalence class. Oppositely labeled points on a shared ray expose precisely the condition that direction-only compression can violate. 

Oehler and Gray\ \citep{OehlerGray1995} incorporate Bayes risk into vector quantization. {\c S}enozan and Soylu\ \citep{Senozan2024} learn entropy-based attribute cut points that accommodate non-monotonic class relationships. Our repair restricts the extra measurement to one common-threshold radius bit and proves its sufficiency, operating interval and error tolerance. Ni et al.\ \citep{Ni2024} study finite-sample classification capacity through repeated affine transformations and normalization, providing a different guarantee from recovery over continuous class regions.

Magnitude retention in attention \citep{Meng2026}, decision-aware logit normalization \citep{Ding2026}, learned graph centering \citep{Scholkemper2025} and adaptive input transformation \citep{September2024} offer mechanisms for controlling retained information. Spatial-sign discriminants use class locations to classify unnormalized observations \citep{Zhuang2026}; their quadratic extension estimates interclass scale \citep{Shen2025}. The present analysis holds the population and direction transformation explicit so that a radial measurement's contribution can be proved separately from learning.

\paragraph{From a chosen projection to every origin.}
Hamsici and Martinez\ \citep{Hamsici2008} optimize projected Bayes error for Gaussian classes with common covariance. Oh\ \citep{Oh2025} addresses merging of nearby classes under linear dimensionality reduction by optimizing generalized means of pairwise scatter ratios, with classification experiments on five datasets. Spherical-model analysis compares decision boundaries with Gaussian approximations \citep{Hamsici2007}, and normalized-feature support vector machines adjust decision offsets \citep{Graf2003}. These approaches clarify the relationship between a chosen representation and its class separation. Our optimization variable is the common origin of a nonlinear radial map.

Hurtado et al.\ \citep{Hurtado2001} study wedge separation of finite planar point sets. Their convex-hull obstruction motivates a continuous-population construction in which shared rays carry opposing probability mass. Orthogonal projections are the geometric limits of radial projections from distant origins \citep{Orponen2019}. Combining these streams identifies the remaining question: when does loss at every finite origin remain bounded away from zero even as the origin escapes to infinity, and what restricted radial measurement restores the decision? The uniform risk limit, signed-density certificates and explicit one-bit decoder provide the corresponding analytical answers.
}

\newcommand{\HalfMoonDiscussion}{%
\section{Discussion}\label{sec:discussion}
The global criterion explains why searching for a better origin can leave an irreducible classification loss. Finite origins merge opposite class preferences along rays, and origins receding in any direction approach a linear projection with its own decision risk. Uniform convergence connects these finite and limiting views into one optimization problem. For connected planar winner regions with mutual convex-hull intrusion, both views lose decisions and a positive global floor follows. This turns decision-preservation criteria \citep{Silva2012,Silva2024,Gyorfi2023} into a geometric test for a whole family of moving-origin representations.

The half-moon results distinguish the guarantee across origins from the risk at a selected origin. The constructive certificates give explicit lower bounds on the best attainable normalization loss, whereas Table~\ref{tab:neural-main} evaluates the specified origins. For zero or negative vertical displacement, a separating strip supplies decision-preserving origins. Every fixed positive displacement in the disjoint regime of Theorem~1 has a positive global floor; the overlapping configuration adds normalization loss to existing class ambiguity. For the original midpoint experiment, $10.835\%$ neural error lies about $0.087$ percentage points above the $10.748\%$ population floor. The decomposition locates nearly all of that observed error in the direction representation.

The one-bit construction identifies exactly which discarded information is needed at $h=0.5$. A near/far indicator selects between the two angular rules in Fig.~\ref{fig:bit-repair}. The zero-bit impossibility result and the explicit one-bit decoder jointly establish the minimum additional radial budget. The Threshold Addendum converts the operating interval into a sharp radius-error guarantee: choosing its midpoint gives the largest tolerance, while the prespecified experimental cut remains unchanged. Within decision-aware compression \citep{OehlerGray1995,Dubois2021}, the construction supplies a simple measurement rule: a coarser radial code preserves the decision while full radius reconstructs the observation. Under the equal-prior reflected sampling laws, radius has the same distribution in both classes, so the bit alone has $50\%$ Bayes error. Its value arises jointly with direction by selecting the appropriate decision rule.

The repair experiment then measures how a trained network uses that sufficient input. Its three errors across 500,000 observations coexist with the decoder's zero observed errors and the proved zero population risk. Together with the original paired-origin study, this supplies a controlled benchmark: determine a representation's attainable risk, hold the learning protocol fixed, and measure learned excess separately. Researchers can thereby distinguish a change in retained decision information from a change in how effectively a learner uses it.
}
